\documentclass[11pt,a4paper]{article}

\usepackage[T1]{fontenc}
\usepackage[utf8]{inputenc}
\usepackage{amsmath,amssymb}
\usepackage{booktabs}
\usepackage{geometry}
\usepackage{hyperref}
\usepackage{xcolor}
\usepackage{listings}
\usepackage{microtype}

\geometry{margin=25mm}
\hypersetup{
  colorlinks=true,
  linkcolor=blue!60!black,
  urlcolor=blue!60!black
}

\lstdefinelanguage{yaml}{
  basicstyle=\ttfamily\small,
  comment=[l]{\#},
  commentstyle=\color{gray},
  morestring=[b]",
  morestring=[b]',
  stringstyle=\color{green!40!black},
  showstringspaces=false,
  breaklines=true,
  columns=fullflexible
}

\title{Analysis of the PCA Scenario QUBO Implementation}
\author{Repository code review}
\date{7 September 2026}

\begin{document}
\maketitle

\section{Executive conclusion}

Yes, the PCA path can encode a categorical choice among up to 21 return bins for
each asset. It imposes one-hot constraints independently for every asset, and the
correlated variant adds pairwise cross-asset compatibility terms. There is,
however, an important distinction between two different uses of the number 21:

\begin{center}
\begin{tabular}{@{}p{0.28\textwidth}p{0.62\textwidth}@{}}
\toprule
Quantity & Meaning \\
\midrule
\texttt{scenariosPerComponents: [21, 5]}
& The outer PCA factor grid contains $21\times5=105$ factor states when two
principal components are used. \\
\texttt{nZBins: 21}
& For each outer PCA state, each asset has up to 21 candidate idiosyncratic
residual-return bins. These are the categorical choices encoded in the QUBO. \\
\bottomrule
\end{tabular}
\end{center}

Consequently, this is not one QUBO containing 105 indivisible market scenarios.
It is normally a stream of 105 separate QUBOs, one for each PCA factor state.
Within each QUBO, the optimizer selects one residual bin per asset. The margin
calculator then takes the greatest decoded portfolio loss across the 105 outer
states.

\section{Scenario construction}

For an outer PCA factor state $s$, the generator first constructs the common
factor return centre
\begin{equation}
  \mathbf{c}_s = \boldsymbol{\mu}_{\mathrm{PCA}}
  + \mathbf{f}_s^{\mathsf T}\mathbf{L},
\end{equation}
where $\mathbf{f}_s$ is the selected principal-component shock and $\mathbf{L}$
maps component shocks back to asset space.

Around this common centre, each asset $i$ receives a discrete candidate set
\begin{equation}
  \mathcal{R}^{(s)}_i
  = \left\{r^{(s)}_{i1},\ldots,r^{(s)}_{iK_i}\right\},
  \qquad K_i \leq 21.
\end{equation}
The candidates are derived from standardized historical residuals and the
asset's residual volatility. The configured residual-sigma range can filter
candidate bins, so ``21'' is a maximum rather than a guarantee that every asset
always has exactly 21 valid choices.

In the reviewed configuration, the relevant grid settings are conceptually:

\begin{lstlisting}[language=yaml]
nComponents: 2
scenariosPerComponents: [21, 5]
nZBins: 21
residualSigmaRange: 5
\end{lstlisting}

With \texttt{nNearest: null}, the implementation may use up to 100 recent
observations from the exponentially weighted fitting window when constructing
the residual grid.

\section{QUBO variables and one-hot constraints}

For every asset $i$ and candidate bin $k$, the encoder creates one binary
variable
\begin{equation}
  x_{ik}\in\{0,1\}.
\end{equation}
The intended categorical constraint is
\begin{equation}
  \sum_{k=1}^{K_i}x_{ik}=1
  \qquad\text{for every asset }i.
\end{equation}

The loss term for portfolio weights $w_i$ is linear in these variables. Using
the sign convention of the implementation, a representative per-state energy is
\begin{equation}
\begin{split}
 E_s(\mathbf{x}) ={}&
   \sum_i\sum_k w_i r^{(s)}_{ik}x_{ik} \\
 &+\lambda_{\mathrm{OH}}
   \sum_i\left(\sum_k x_{ik}-1\right)^2 \\
 &+\lambda_{\mathrm{compat}}
   \sum_{(i,j)\in\mathcal{E}}\sum_{k,\ell}
   C^{(s)}_{ij,k\ell}x_{ik}x_{j\ell}.
\end{split}
\end{equation}

Expanding the one-hot penalty gives the required linear and quadratic biases:
\begin{equation}
 \left(\sum_k x_{ik}-1\right)^2
 =-\sum_k x_{ik}+2\sum_{k<\ell}x_{ik}x_{i\ell}+1,
\end{equation}
because $x_{ik}^2=x_{ik}$ for binary variables. Thus each asset constitutes its
own one-hot group, and choosing zero or multiple bins is penalized.

The solver path first prefers the lowest-energy feasible one-hot sample. If a
solver supplies no feasible candidate, the framework projects the returned
sample into the categorical groups and applies deterministic local improvement
before decoding. This post-processing is a defensive mechanism; it does not
replace the one-hot penalty encoded in the QUBO.

\section{Cross-asset compatibility}

Cross-asset compatibility is present only when the
\texttt{correlated\_returns\_vola\_grid} generator is used. For the plain
\texttt{returns\_vola\_grid} generator, the BQM visitor explicitly sets the
compatibility penalty to zero.

For the correlated generator, asset pairs are built from the undirected union
of directed top-$k$ absolute-correlation nominations. With \texttt{topK: 5}, an
asset nominates its five strongest correlation neighbours, and an edge is kept
if either endpoint nominates the other.

For an edge $(i,j)$ and standardized candidate residuals $z_{ik}$ and
$z_{j\ell}$, the compatibility cost is based on a bivariate Gaussian
Mahalanobis form:
\begin{equation}
  C_{ij,k\ell}
  =\frac{z_{ik}^2-2\rho_{ij}z_{ik}z_{j\ell}+z_{j\ell}^2}
  {1-\rho_{ij}^2}.
\end{equation}
This cost is multiplied by $\lambda_{\mathrm{compat}}$ and inserted as a
quadratic coupling between $x_{ik}$ and $x_{j\ell}$.

The compatibility term is soft, not a hard prohibition. It discourages
statistically implausible combinations, but a sufficiently adverse portfolio
loss can outweigh it depending on the penalty scale.

\section{What remains joint and what is recombined}

The PCA construction has two dependence layers:

\begin{enumerate}
  \item The outer PCA point is a joint common-factor shock applied to every
  asset simultaneously.
  \item Conditional on that PCA point, the QUBO selects an idiosyncratic
  residual bin separately for each asset. In the correlated generator, pairwise
  compatibility costs regularize these selections.
\end{enumerate}

Therefore the PCA QUBO does perform cross-asset recombination of residual bins.
It does not treat each historical or generated row as an indivisible joint
scenario. This is intentional for the PCA grid model.

That behavior differs fundamentally from FHS--EVT. An FHS--EVT risk state
contains one synchronized return for every instrument, and the entire row is an
indivisible joint scenario. A correct FHS--EVT evaluator must never take one
asset's shock from one row and another asset's shock from a different row.

\begin{center}
\begin{tabular}{@{}p{0.26\textwidth}p{0.31\textwidth}p{0.33\textwidth}@{}}
\toprule
Property & PCA return/volatility grid & FHS--EVT \\
\midrule
Outer state
& One common PCA-factor point
& One complete synchronized joint scenario \\
Inner asset choice
& One residual bin per asset
& None \\
Cross-asset handling
& PCA factor dependence plus optional soft pairwise residual compatibility
& Dependence already embedded in the scenario row \\
Correct optimization unit
& Categorical asset-bin assignment conditional on a PCA point
& Whole scenario row \\
\bottomrule
\end{tabular}
\end{center}

\section{Configuration-specific answer}

The standard synthetic 200-asset BQM configuration uses the correlated PCA
generator and therefore has both one-hot constraints and cross-asset
compatibility. Its key settings include:

\begin{lstlisting}[language=yaml]
riskStateGenerator:
  type: correlated_returns_vola_grid
  nComponents: 2
  scenariosPerComponents: [21, 5]
  nZBins: 21
  topK: 5

marginCalculator:
  type: bqm
  lambdaOneHot: 1.0
  lambdaCompat: 0.1
\end{lstlisting}

By contrast, the latest Yahoo 200-asset PCA-versus-FHS--EVT comparison used
\texttt{returns\_vola\_grid} with the greedy calculator. That particular PCA
comparator did not invoke a QUBO and did not apply cross-asset compatibility.

\section{Direct answer}

For the PCA BQM configuration, the answer is yes: the QUBO chooses one of up to
21 residual bins for each asset, enforces one-hot selection for every asset, and
the correlated variant incorporates pairwise cross-asset compatibility. It does
this separately for each of the 105 outer PCA factor points and takes the worst
decoded portfolio result across those points.

The 105 PCA points and the 21 per-asset QUBO bins are distinct dimensions and
should not be conflated. Nor should the PCA construction be imposed on FHS--EVT,
whose output rows must remain indivisible joint scenarios.

\section{Implementation references}

The conclusions above were traced through the following repository files:

\begin{itemize}
  \item \path{src/risk_state_generator/returns_vola_grid_risk_state_generator.py}
  \item \path{src/risk_state_generator/correlated_returns_vola_grid_risk_state_generator.py}
  \item \path{src/risk_state_generator/portfolio_risk_state_bqm_visitor.py}
  \item \path{src/margin_calculator/bqm_margin_calculator.py}
  \item \path{config/backtests/assets_00200.yaml}
  \item \path{config/backtests/yahoo_00200_pca_four_year_full.yaml}
  \item \path{config/backtests/yahoo_00200_fhs_evt_ar2_tail075_expanding_semiannual_full_nodes2_xi20.yaml}
\end{itemize}

This document records a read-only analysis. No source-code alterations were made
as part of the underlying QUBO review.

\end{document}
